
\documentclass{article}
\usepackage{colortbl}
\usepackage{makecell}
\usepackage{multirow}
\usepackage{supertabular}

\begin{document}

\newcounter{utterance}

\centering \large Interaction Transcript for game `cladder', experiment `full\_v1.5\_default', episode 8797 with Qwen3.5{-}4B.
\vspace{24pt}

{ \footnotesize  \setcounter{utterance}{1}
\setlength{\tabcolsep}{0pt}
\begin{supertabular}{c@{$\;$}|p{.15\linewidth}@{}p{.15\linewidth}p{.15\linewidth}p{.15\linewidth}p{.15\linewidth}p{.15\linewidth}}
    \# & \multicolumn{2}{c}{Player} && \multicolumn{2}{c}{Game Master} \\
    \hline

    \theutterance \stepcounter{utterance}  
    & & & \multicolumn{4}{p{0.6\linewidth}}{
        \cellcolor[rgb]{0.9,0.9,0.9}{
            \makecell[{{p{\linewidth}}}]{
                \texttt{\tiny{[P1$\langle$GM]}}
                \texttt{You are an expert in causal inference. The following question is not a typical commonsense query, but rather a meticulously designed question created by a professor specializing in causal inference, intended to assess the students' mastery of the course content.} \\
\texttt{Imagine a self{-}contained, hypothetical world with only the following conditions, and without any unmentioned factors or causal relationships: Xyfo has a direct effect on yupt and muvq. Yupt has a direct effect on zupj. Zupj has a direct effect on muvq. Xyfo is unobserved. For those who are not yupt and are not zupj, the probability of muvq is 90\%. For those who are not yupt and are zupj, the probability of muvq is 54\%. For those who are yupt and are not zupj, the probability of muvq is 96\%. For those who are yupt and are zupj, the probability of muvq is 57\%. For those who are not yupt, the probability of zupj is 51\%. For those who are yupt, the probability of zupj is 34\%. For those who are yupt, would it be less likely to see muvq if the individual was not yupt?} \\
\texttt{Start your answer with 'Yes' or 'No', followed by additional reasoning or evidence to support your explanation.} \\
            }
        }
    }
     \\ \\

    \theutterance \stepcounter{utterance}  
    & & \multicolumn{4}{p{0.6\linewidth}}{
        \cellcolor[rgb]{0.95,0.95,0.95}{
            \makecell[{{p{\linewidth}}}]{
                \texttt{\tiny{[P1$\leftrightarrow$P]}}
                \texttt{{[}DRAFT{]}} \\
\texttt{No} \\
\\ 
\texttt{**Reasoning:**} \\
\\ 
\texttt{The question asks whether, for individuals who are **yupt**, it would be **less likely** to see **muvq** if they were **not yupt**. This is a counterfactual question asking for the comparison between the potential outcome \$Y\_\{muvq\}(1)\$ (outcome if yupt) and \$Y\_\{muvq\}(0)\$ (outcome if not yupt) for the subpopulation defined by \$Y\_\{yupt\}=1\$.} \\
\\ 
\texttt{1.  **Analyze the Counterfactual Logic**:} \\
\texttt{    The phrase "For those who are yupt, would it be less likely to see muvq if the individual was not yupt?" translates mathematically to asking if the probability of \$muvq\$ under the treatment \$Y\_\{yupt\}=1\$ is lower than the probability of \$muvq\$ under the control \$Y\_\{yupt\}=0\$.} \\
\texttt{    In causal terms, we are comparing the Average Treatment Effect (ATE) or the difference in potential outcomes:} \\
\texttt{    \$\$P(Y\_\{muvq\}=1 \textbackslash{}mid Y\_\{yupt\}=1) \textbackslash{}quad \textbackslash{}text\{vs.\} \textbackslash{}quad P(Y\_\{muvq\}=1 \textbackslash{}mid Y\_\{yupt\}=0)\$\$} \\
\texttt{    If the first value is smaller than the second, the answer is "Yes". If the first value is larger, the answer is "No".} \\
\\ 
\texttt{2.  **Evaluate the Provided Probabilities**:} \\
\texttt{    The problem provides specific probabilities for the outcome \$muvq\$ based on the status of \$yupt\$ and \$zupj\$.} \\
\texttt{    *   **Group A (yupt = 1)**:} \\
\texttt{        *   If \$zupj=0\$: \$P(muvq) = 96\textbackslash{}\%\$} \\
\texttt{        *   If \$zupj=1\$: \$P(muvq) = 57\textbackslash{}\%\$} \\
\texttt{        *   To find the overall probability for those who are yupt, we must weight these by the probability of \$zupj\$ given \$yupt\$.} \\
\texttt{        *   Given: \$P(zupj=1 \textbackslash{}mid yupt=1) = 34\textbackslash{}\%\$ (so \$P(zupj=0 \textbackslash{}mid yupt=1) = 66\textbackslash{}\%\$).} \\
\texttt{        *   \$P(muvq \textbackslash{}mid yupt=1) = (0.96 \textbackslash{}times 0.66) + (0.57 \textbackslash{}times 0.34) = 0.6336 + 0.1938 = 0.8274\$ (or **82.74\%**).} \\
\\ 
\texttt{    *   **Group B (yupt = 0)**:} \\
\texttt{        *   If \$zupj=0\$: \$P(muvq) =} \\
            }
        }
    }
    & \\ \\

    \theutterance \stepcounter{utterance}  
    & & \multicolumn{4}{p{0.6\linewidth}}{
        \cellcolor[rgb]{0.95,0.95,0.95}{
            \makecell[{{p{\linewidth}}}]{
                \texttt{\tiny{[P1$\leftrightarrow$P]}}
                \texttt{{[}SELF{-}FEEDBACK{]}} \\
\texttt{FORMAT: problem} \\
\texttt{RULES: problem} \\
\texttt{STRATEGY: problem} \\
\texttt{REASON: The answer fails to complete the calculation for Group B, does not provide a final conclusion, and ignores the required "Yes" or "No" prefix format.} \\
            }
        }
    }
    & \\ \\

    \theutterance \stepcounter{utterance}  
    & \multicolumn{4}{p{0.6\linewidth}}{
        \cellcolor[rgb]{0.561,0.737,0.561}{
            \makecell[{{p{\linewidth}}}]{
                \texttt{\tiny{[P1$\rangle$GM]}}
                \texttt{No} \\
\\ 
\texttt{**Reasoning:**} \\
\\ 
\texttt{The question asks whether, for individuals who are **yupt**, it would be **less likely** to see **muvq** if they were **not yupt**. This is a counterfactual question asking for the comparison between the probability of \$muvq\$ under the treatment \$Y\_\{yupt\}=1\$ and the probability of \$muvq\$ under the control \$Y\_\{yupt\}=0\$.} \\
\\ 
\texttt{1.  **Analyze the Counterfactual Logic**:} \\
\texttt{    We need to compare \$P(muvq \textbackslash{}mid yupt=1)\$ with \$P(muvq \textbackslash{}mid yupt=0)\$. If the probability in the first group is lower than in the second, the answer is "Yes"; otherwise, it is "No".} \\
\\ 
\texttt{2.  **Evaluate the Provided Probabilities**:} \\
\texttt{    *   **Group A (yupt = 1)**:} \\
\texttt{        *   If \$zupj=0\$: \$P(muvq) = 96\textbackslash{}\%\$} \\
\texttt{        *   If \$zupj=1\$: \$P(muvq) = 57\textbackslash{}\%\$} \\
\texttt{        *   Given: \$P(zupj=1 \textbackslash{}mid yupt=1) = 34\textbackslash{}\%\$ (so \$P(zupj=0 \textbackslash{}mid yupt=1) = 66\textbackslash{}\%\$).} \\
\texttt{        *   \$P(muvq \textbackslash{}mid yupt=1) = (0.96 \textbackslash{}times 0.66) + (0.57 \textbackslash{}times 0.34) = 0.6336 + 0.1938 = 0.8274\$ (or **82.74\%**).} \\
\\ 
\texttt{    *   **Group B (yupt = 0)**:} \\
\texttt{        *   If \$zupj=0\$: \$P(muvq) = 90\textbackslash{}\%\$} \\
\texttt{        *   If \$zupj=1\$: \$P(muvq) = 54\textbackslash{}\%\$} \\
\texttt{        *   Given: \$P(zupj=1 \textbackslash{}mid yupt=0) = 51\textbackslash{}\%\$ (so \$P(zupj=0 \textbackslash{}mid yupt=0) = 49\textbackslash{}\%\$).} \\
\texttt{        *   \$P(muvq \textbackslash{}mid yupt=0) = (0.90 \textbackslash{}times 0.49) + (0.54 \textbackslash{}times 0.51) = 0.4410 + 0.2754 = 0.7164\$ (or **71.64\%**).} \\
\\ 
\texttt{3.  **Compare and Conclude**:} \\
\texttt{    Comparing the two probabilities:} \\
\texttt{    \$P(muvq \textbackslash{}mid yupt=1) =} \\
            }
        }
    }
    & & \\ \\

    \theutterance \stepcounter{utterance}  
    & & & \multicolumn{2}{p{0.3\linewidth}}{
        \cellcolor[rgb]{0.9,0.9,0.9}{
            \makecell[{{p{\linewidth}}}]{
                \texttt{\tiny{[GM$|$GM]}}
                \texttt{No} \\
            }
        }
    }
    & & \\ \\

    \theutterance \stepcounter{utterance}  
    & & & \multicolumn{2}{p{0.3\linewidth}}{
        \cellcolor[rgb]{0.9,0.9,0.9}{
            \makecell[{{p{\linewidth}}}]{
                \texttt{\tiny{[GM$|$GM]}}
                \texttt{yes} \\
            }
        }
    }
    & & \\ \\

    \theutterance \stepcounter{utterance}  
    & & & \multicolumn{2}{p{0.3\linewidth}}{
        \cellcolor[rgb]{0.9,0.9,0.9}{
            \makecell[{{p{\linewidth}}}]{
                \texttt{\tiny{[GM$|$GM]}}
                \texttt{game\_result = LOSE} \\
            }
        }
    }
    & & \\ \\

\end{supertabular}
}

\end{document}
